\begin{center}
\large{\MakeUppercase{\textbf{Анотація}}}
\end{center}				
%\addcontentsline{toc}{chapter}{ВИСНОВКИ}	% Добавляем его в оглавление


\newcommand{\actuality}{\textbf{Актуальність теми.}}
\newcommand{\aim}{\textbf{Метою}}
\newcommand{\tasks}{\textbf{задачі}}
\newcommand{\contribution}{\textbf{Особистий вклад.}}
\newcommand{\keyword}{\textbf{Ключові слова:}}

\textbf{Об'єм і структура проєкту.} Проєкт складається з вступу, чотирьох розділів, висновків, переліку умовних
скорочень та списку літератури.

{\actuality} Роботу присвячено розробці програмно-алгоритмічного комплексу для автоматизованого відстеження волейбольного
м'яча та розрахунку його кінематичних параметрів(швидкості, 3D-траєкторії) на основі монокулярного відеопотоку.

В процесі виконання роботи проаналізовано існуючі комерційні системи спортивної аналітики та виявлено потребу в
доступних мобільних рішеннях.
Розроблено архітектуру системи, що базується на моделі глибокого навчання YOLO для
високоточної об'єктної детекції в умовах розмиття рухом.\ Для переходу від двовимірних піксельних
координат до фізичних метричних величин імплементовано модель камери-обскури з використанням алгоритму просторового
калібрування PnP та методу зворотного проектування.

Для вирішення проблеми високочастотного апаратного шуму та компенсації нелінійних аеродинамічних
ефектів застосовано математичний апарат оптимальної фільтрації — сигма-точковий фільтр Калмана.
З метою точного відстеження різких змін швидкості під час атакуючих ударів розроблено гібридну
архітектуру на базі взаємодіючих множинних моделей.

Експериментальне тестування підтвердило здатність системи реконструювати тривимірну траєкторію та фіксувати пікову
швидкість удару з використанням однієї камери стандартної частоти кадрів (50-60 FPS). Результати роботи можуть бути
використані тренерськими штабами для тактико-технічного аналізу та оптимізації тренувального процесу у волейболі.

%
\keyword\ \textit{модель глибокого навчання YOLO, алгоритм просторового калібрування PnP, метод зворотного проектування,
    нелінійні аеродинамічні ефекти, сигма-точковий фільтр Калмана, взаємодіючі множинні моделі, тривімірна
    траєкторія, пікова швидкість удару.}


\newpage

\begin{center}
\large{\MakeUppercase{\textbf{Annotation}}}
\end{center}

The thesis is devoted to the development of a software and algorithmic complex for automated volleyball tracking and
calculation of its kinematic parameters (speed, 3D trajectory) based on a monocular video stream.

During the research, existing commercial sports analytics systems were analyzed, revealing a strong demand for
accessible mobile solutions.\ The system architecture is based on the YOLO deep learning model for high-precision object
detection under severe motion blur conditions.\ To transition from 2D pixel coordinates to physical metric values, a
pinhole camera model was implemented using the Perspective-n-Point(PnP) spatial calibration algorithm and Ray Casting
method.

To solve the problem of high-frequency hardware noise and compensate for nonlinear aerodynamic
effects, the mathematical apparatus of optimal filtering was applied, specifically the Unscented Kalman Filter.
In order to accurately track abrupt speed changes during spikes, a hybrid architecture based on Interacting Multiple
Models was designed.

Experimental testing confirmed the system's ability to reconstruct the 3D trajectory and capture the peak spike speed
using a single standard frame rate camera (50-60 FPS). The results of the work can be utilized by coaching staffs for
tactical and technical analysis and the optimization of the training process in volleyball.

\textbf{Keywords:} \textit{YOLO deep learning model, PnP spatial calibration algorithm, Ray Casting
    method, nonlinear aerodynamic effects, Unscented Kalman Filter, Interacting Multiple Models three-dimensional,
    trajectory, peak impact velocity.}

\newpage